%==========================================================
%  Conclusion.tex  —  Conclusion générale
%==========================================================
\chapter*{Conclusion générale}
\addcontentsline{toc}{chapter}{Conclusion générale}
\markboth{Conclusion générale}{}

\thispagestyle{fancy}

\definecolor{tekupBlue}{RGB}{0,70,127}
\definecolor{lightBlue}{RGB}{230,240,250}

%----------------------------------------------------------
% Bilan du projet
%----------------------------------------------------------
Ce projet de fin d'études, réalisé au sein de \textbf{3LM Solution}, a abouti à
la conception et au développement d'une plateforme CRM intelligente pour
centres d'appels, intégrant des fonctionnalités avancées d'intelligence
artificielle pour l'analyse et le scoring des appels.

Les principaux résultats obtenus couvrent l'ensemble du périmètre fonctionnel
initialement défini~: un module d'authentification et de gestion des
utilisateurs avec un contrôle d'accès basé sur les permissions, une gestion
complète des leads avec qualification et suivi, un système de rendez-vous et
de campagnes, un moteur de transcription et d'analyse des appels basé sur
Whisper et un LLM local (gemma3:4b), un module d'évaluation de la qualité
avec scoring automatique et manuel, ainsi que des fonctionnalités de gestion
de présence, de paie et de reporting.

%----------------------------------------------------------
% Apports techniques et méthodologiques
%----------------------------------------------------------
D'un point de vue technique, ce projet m'a permis de maîtriser l'ensemble de
la stack \textbf{.NET 8 / ASP.NET Core} avec une architecture clean
architecture respectant les principes du Domain-Driven Design, couplée à un
frontend moderne en \textbf{React 18 / TypeScript}. L'intégration de modèles
d'IA via \textbf{Ollama} a constitué un défi technique particulièrement
enrichissant, notamment pour l'orchestration des appels au LLM avec gestion
de timeout et de température, ainsi que pour le pipeline de transcription
audio.

La méthodologie \textbf{Scrum} adoptée tout au long du projet a permis une
livraison progressive et itérative des fonctionnalités, avec une priorisation
continue des user stories et une adaptation constante aux retours de
l'encadrant professionnel.

%----------------------------------------------------------
% Limites et perspectives
%----------------------------------------------------------
Malgré les résultats satisfaisants obtenus, plusieurs pistes d'amélioration
méritent d'être explorées~:
\begin{itemize}
    \item \textbf{Analyse en temps réel}~: le déploiement de l'analyse
    sentimentale en continu pendant les appels, aujourd'hui désactivée par
    défaut, pourrait être activé pour fournir des feedbacks instantanés aux
    agents.
    \item \textbf{Diarisation}~: l'identification du locuteur (agent vs client)
    dans les transcriptions pourrait être améliorée avec des modèles dédiés.
    \item \textbf{Dashboard prédictif}~: l'ajout de modèles prédictifs pour
    anticiper le taux de conversion des leads et les risques d'attrition.
    \item \textbf{Déploiement cloud}~: le déploiement sur une infrastructure
    cloud scalable (Azure, AWS) avec mise en place d'un load balancing pour
    le service d'IA.
\end{itemize}

\vspace{0.5cm}

\begin{center}
    \begin{tcolorbox}[
        colback=lightBlue,
        colframe=tekupBlue,
        boxrule=0.8pt,
        arc=6pt,
        width=0.85\linewidth,
        halign=center
    ]
        \itshape
        «~Le CRM intelligent n'est pas seulement un outil de gestion de la relation
        client~; c'est un levier stratégique pour transformer chaque interaction
        en opportunité de valeur.~»
    \end{tcolorbox}
\end{center}
